% ==============================================================================
% SLIDES - SEMANA 07: GESTÃO DE DESEMPENHO EM TI, KPIS & MÉTRICAS
% ==============================================================================

% ------------------------------------------------------------------------------
% 1. SLIDE DE CAPA
% ------------------------------------------------------------------------------
\senaicover
  {Gestão de Desempenho em TI}
  {Métricas, KPIs do Negócio vs. Operacionais e Dashboard Executivo}
  {Unidade Curricular: Governança de TI}
  {Prof. André Souza}

% ------------------------------------------------------------------------------
% 2. VISÃO GERAL / AGENDA
% ------------------------------------------------------------------------------
\begin{senaiframe}{Visão Geral \& Agenda}{Os 4 Eixos da Aula}

  \begin{columns}[t]
    \begin{column}{0.48\textwidth}
      \begin{tcolorbox}[senaibluecard, title={1. Hierarquia da Medição}]
        \begin{itemize}
          \item Dados vs. Métricas vs. KPIs.
          \item A regra de Peter Drucker na TI.
        \end{itemize}
      \end{tcolorbox}
      
      \vspace{4pt}
      
      \begin{tcolorbox}[senaiambercard, title={3. Linha de Base \& Metas}]
        \begin{itemize}
          \item Estabelecendo a Baseline atual.
          \item Definição de Metas SMART e frequência.
        \end{itemize}
      \end{tcolorbox}
    \end{column}

    \begin{column}{0.48\textwidth}
      \begin{tcolorbox}[senairedcard, title={2. Categorias de KPIs}]
        \begin{itemize}
          \item KPIs do Negócio (ROI, Receita Digital, CSAT).
          \item KPIs Operacionais (MTTR, MTBF, Availability).
        \end{itemize}
      \end{tcolorbox}
      
      \vspace{4pt}
      
      \begin{tcolorbox}[senaigreencard, title={4. Capítulo 5 do PDGTI}]
        \begin{itemize}
          \item Sistema de Indicadores de Desempenho.
          \item Mock-up do Dashboard Executivo de TI.
        \end{itemize}
      \end{tcolorbox}
    \end{column}
  \end{columns}

\end{senaiframe}

% ------------------------------------------------------------------------------
% 3. KPIS DO NEGÓCIO VS OPERACIONAIS
% ------------------------------------------------------------------------------
\begin{senaiframe}{Módulo 1 • Indicadores}{KPIs Executivos vs. KPIs Operacionais}

  \begin{columns}[t]
    \begin{column}{0.48\textwidth}
      \begin{tcolorbox}[senaibluecard, title={KPIs do Negócio (Estratégico)}]
        \begin{itemize}
          \item \textbf{ROI de TI:} Retorno financeiro sobre investimentos em projetos.
          \item \textbf{CSAT / NPS de TI:} Nível de satisfação dos usuários internos.
          \item \textbf{Time-to-Market:} Tempo para lançar novos produtos digitais.
        \end{itemize}
      \end{tcolorbox}
    \end{column}

    \begin{column}{0.48\textwidth}
      \begin{tcolorbox}[senaicard, title={KPIs Operacionais (Técnico)}]
        \begin{itemize}
          \item \textbf{Disponibilidade / Uptime (%):} Tempo de atividade dos sistemas críticos.
          \item \textbf{MTTR:} Tempo Médio para Reparo de incidentes.
          \item \textbf{MTBF:} Tempo Médio Entre Falhas consecutivas.
        \end{itemize}
      \end{tcolorbox}
    \end{column}
  \end{columns}

\end{senaiframe}

% ------------------------------------------------------------------------------
% 4. METAS SMART & DASHBOARD
% ------------------------------------------------------------------------------
\begin{senaiframe}{Módulo 2 • Metas \& Painel}{Estrutura de um KPI de Sucesso}

  \begin{columns}[t]
    \begin{column}{0.48\textwidth}
      \begin{tcolorbox}[senaiambercard, title={Metas SMART}]
        \begin{itemize}
          \item \textbf{Específica:} Específica e sem ambiguidade.
          \item \textbf{Mensurável:} Quantificável numericamente.
          \item \textbf{Atingível \& Relevante:} Desafiadora, mas viável.
          \item \textbf{Temporal:} Prazo limite claro para alcance.
        \end{itemize}
      \end{tcolorbox}
    \end{column}

    \begin{column}{0.48\textwidth}
      \begin{tcolorbox}[senaigreencard, title={Dashboard Executivo}]
        \begin{itemize}
          \item Coleta automatizada sem sobrecarga manual.
          \item Apresentação visual intuitiva para o Comitê de TI (estilo Farol/Semáforo: Verde, Amarelo, Vermelho).
        \end{itemize}
      \end{tcolorbox}
    \end{column}
  \end{columns}

\end{senaiframe}

% ------------------------------------------------------------------------------
% 5. REFERÊNCIAS BIBLIOGRÁFICAS
% ------------------------------------------------------------------------------
\begin{senaiframe}{Referências Bibliográficas}{Fontes \& Leitura Recomendada}

  \begin{tcolorbox}[senaicard, title={Obras de Referência}]
    \begin{itemize}
      \item \textbf{FERNANDES, Aguinaldo Aragon; ABREU, Vladimir Ferraz.} *Implantando a governança de TI: da estratégia à gestão de processos e serviços.* 4. ed. Rio de Janeiro: Brasport, 2014.
      \item \textbf{BITTENCOURT, Carlos Magno Andriolli.} *Governança corporativa e compliance: planejamento e gestão estratégica.* Curitiba: Contentus, 2020.
    \end{itemize}
  \end{tcolorbox}

\end{senaiframe}

% ------------------------------------------------------------------------------
% 6. APLICAÇÃO NO PDGTI
% ------------------------------------------------------------------------------
\begin{senaiframe}{Aplicação no PDGTI}{Desenvolvimento do Capítulo 5}

  \vspace{0.2cm}
  \begin{tcolorbox}[senaigreencard, title={Capítulo 5 – Gestão de Desempenho e Indicadores}]
    \begin{itemize}
      \item \textbf{5.1 KPIs do Negócio:} Definir no mínimo 3 KPIs com impacto estratégico.
      \item \textbf{5.2 KPIs Operacionais de TI:} Definir no mínimo 4 KPIs técnicos (Uptime, MTTR, SLA).
      \item \textbf{5.3 Linha de Base e Metas:} Estabelecer baseline, meta SMART, responsável e apuração.
      \item \textbf{5.4 Dashboard Executivo:} Mock-up visual do painel a ser apresentado ao Comitê.
    \end{itemize}
  \end{tcolorbox}

\end{senaiframe}
